% Informe — plantilla relajada con portada, cajas destacadas, diagrama y código.
% Tectonic es XeTeX: las fuentes ya son UTF-8, sin paquete inputenc.
% spell-whitelist.txt cubre los valores de ejemplo en inglés y las claves de los diagramas.
% Reemplace cada bloque de ejemplo y cada comentario con % por su propio contenido.
\documentclass[11pt,a4paper]{article}

% --- Preámbulo (solo paquetes del bundle de Tectonic; sin inputenc/listings/minted) ---
\usepackage[spanish,es-nodecimaldot]{babel}
\usepackage[T1]{fontenc}
\usepackage{lmodern}
\usepackage{amssymb}
\usepackage[a4paper,margin=2.5cm,headheight=14pt,headsep=12pt,footskip=14pt]{geometry}
\usepackage{xcolor}
\usepackage{hyperref}
\usepackage{enumitem}
\usepackage{fancyhdr}
\usepackage{titlesec}
\usepackage{tcolorbox}
\usepackage{booktabs}
\usepackage{float}

% --- Acento único: un solo color para títulos, enlaces, cajas y regla de cabecera ---
\definecolor{acento}{HTML}{1B5FAA}
\definecolor{aviso}{HTML}{B45309}

\hypersetup{colorlinks=true, linkcolor=acento, urlcolor=acento, citecolor=acento,
  breaklinks=true, pdftitle={{{title}}}, pdfauthor={{{author}}}}

% --- Cabecera: curso a la izquierda, título a la derecha, página centrada abajo ---
% % Mantenga {{title}} y {{course}} cortos para que la cabecera no se desborde.
\pagestyle{fancy}
\fancyhf{}
\lhead{{{course}}}
\rhead{{{title}}}
\cfoot{\thepage}
\renewcommand{\headrulewidth}{0.8pt}
\renewcommand{\headrule}{\hbox to\headwidth{\color{acento}\leaders\hrule height \headrulewidth\hfill}}

% --- Títulos de sección en acento ---
\titleformat{\section}{\large\bfseries\color{acento}}{\thesection.}{0.5em}{}
\titleformat{\subsection}{\normalsize\bfseries\color{acento!80!black}}{\thesubsection.}{0.5em}{}
\titlespacing*{\section}{0pt}{8pt}{4pt}
\titlespacing*{\subsection}{0pt}{6pt}{3pt}

\setlength{\parskip}{3pt}
\setlist[itemize]{leftmargin=1.4em, itemsep=1pt, topsep=2pt}
\setlist[enumerate]{leftmargin=1.4em, itemsep=1pt, topsep=2pt}

% --- Cajas reutilizables (cada una se usa una vez en el ejemplo) ---
\newtcolorbox{nota}{colback=gray!4, colframe=gray!45, fonttitle=\bfseries,
  left=4pt, right=4pt, top=3pt, bottom=3pt}
\newtcolorbox{clave}{colback=acento!6, colframe=acento!65!black, fonttitle=\bfseries,
  title=Idea clave, left=4pt, right=4pt, top=3pt, bottom=3pt}
\newtcolorbox{cuidado}{colback=aviso!6, colframe=aviso!70!black, fonttitle=\bfseries,
  title=Atención, left=4pt, right=4pt, top=3pt, bottom=3pt}

\begin{document}

  % --- Portada (debe ocupar exactamente 1 página: no agregar texto aquí) ---
  \begin{titlepage}
    \centering
    % % Reemplace el antetítulo (por ejemplo, el tipo de entrega).
    {\scshape\large {{kicker}} \par}
    \vspace{1.0cm}
    % % Reemplace por el título del documento.
    {\huge\bfseries {{title}} \par}
    \vspace{0.4cm}
    % % Reemplace por una frase que describa el documento.
    {\Large {{subtitle}} \par}
    \vspace{1.4cm}
    \begin{tcolorbox}[colback=acento!4, colframe=acento!45!black, width=0.88\textwidth]
      \begin{tabular}{@{}ll@{}}
        \textbf{Autor/a:} & {{author}} \\
        \textbf{Asignatura o proyecto:} & {{course}} \\
        \textbf{Fecha:} & {{date}} \\
        \textbf{Repositorio:} & {\small\url{{{repository}}}} \\
      \end{tabular}
    \end{tcolorbox}
    \vfill
  \end{titlepage}

  % % Resumen: de tres a cuatro líneas con el objetivo y el resultado principal. Va dentro de `clave`.
  \section{Resumen}
  \begin{clave}
    % % Reemplace por el mensaje principal del documento.
    Este informe resume el avance del trabajo, la solución propuesta y los
    resultados principales en una lectura rápida.
  \end{clave}

  % % Contexto: un párrafo corto (qué se pidió, qué restricción importa, qué se entrega) y la caja `nota`.
  \section{Contexto}
  % % Reemplace por el contexto de su entrega o nota técnica.
  El proyecto necesita una actualización breve del estado actual: qué se pidió,
  qué restricción importa y qué se entrega esta vez.
  \begin{nota}
    % % Reemplace por el alcance o la suposición que el lector debe conocer.
    Alcance de esta entrega: solo lo completado hasta la fecha indicada en la portada.
  \end{nota}

  % % Solución: un párrafo, el diagrama y el bloque de código (máximo 15 líneas cortas).
  \section{Solución}
  % % Reemplace por los pasos de su solución.
  El flujo general tiene tres pasos: preparar los datos, aplicar la regla
  principal y presentar el resultado.
  % % Reemplace por los pasos de su propio flujo; mantenga el estilo y el título.
  \begin{mermaid}[style=editorial, width=0.85\linewidth, caption={Flujo general de la solución}, pos=H]
    flowchart LR
    A[Preparar datos] --> B[Aplicar regla] --> C[Presentar resultado]
  \end{mermaid}
  % % Reemplace por su ejemplo mínimo (máximo 15 líneas de menos de 90 columnas).
  \begin{code}[lang=python, numbers=true]
def resumir(items):
    total = sum(items)
    n = len(items)
    promedio = total / n if n else 0
    return total, promedio
  \end{code}

  % % Resultados: un párrafo, la tabla con título y la caja `cuidado`.
  \section{Resultados}
  % % Reemplace por los valores de su caso.
  La tabla resume tres mediciones del ejemplo con valores ilustrativos.
  \begin{table}[H]
    \centering
    \caption{Mediciones ilustrativas del ejemplo}
    \begin{tabular}{@{}lcc@{}}
      \toprule
      Caso & Antes & Después \\
      \midrule
      A & 12 & 8 \\
      B & 20 & 11 \\
      C & 15 & 9 \\
      \bottomrule
    \end{tabular}
  \end{table}
  \begin{cuidado}
    % % Reemplace por la limitación que el lector no debe pasar por alto.
    Los valores son ilustrativos; verifique las cifras con los datos finales.
  \end{cuidado}

  % % Conclusiones: dos o tres frases de cierre y la lista de pendientes.
  \section{Conclusiones y siguientes pasos}
  % % Reemplace por su cierre y sus pendientes.
  El trabajo cumple el objetivo planteado y deja dos pasos claros para continuar.
  % % Reemplace por sus pendientes; la casilla marca cada tarea abierta.
  \begin{itemize}[label=$\square$]
    \item Revisar el borrador con el equipo.
    \item Incorporar los datos finales.
    \item Publicar la siguiente versión.
  \end{itemize}

  % % Referencias: dos entradas genéricas de ejemplo, sin archivo .bib.
  % % `thebibliography` ya genera el título «Referencias»; no agregue \section* aquí.
  \begin{thebibliography}{9}
    \bibitem{guia} Autor de ejemplo, \textit{Guía breve de informes técnicos}, Editorial de ejemplo, 2024.
    \bibitem{manual} Autora de ejemplo, \textit{Manual de escritura clara}, Editorial de ejemplo, 2023.
  \end{thebibliography}

\end{document}
